\chapter{Tổ chức thực hiện và kết luận}
\label{chap:ket-luan}

\section{Thành viên và vai trò Scrum}
Nhóm 7 có năm thành viên. Vai trò được đối chiếu từ tài liệu phân công
của nhóm~\cite{roles,assignment}; họ tên và MSSV được nhóm cung cấp.
Phan Anh Khang kiêm Product Owner, Scrum Master và Tech Lead/Schema
Owner; cả năm thành viên là Development Team và cùng chịu trách nhiệm
về chất lượng sản phẩm.

\begin{table}[H]
\centering\small
\caption{Vai trò Scrum và trách nhiệm điều phối}
\label{tab:scrumroles}
\begin{tabularx}{\textwidth}{@{}L{0.22\textwidth}L{0.23\textwidth}Y@{}}
\toprule\textbf{Vai trò} & \textbf{Phụ trách} & \textbf{Trách nhiệm} \\\midrule
Product Owner & Phan Anh Khang & Ưu tiên backlog, làm rõ AC và điều phối nghiệm thu Increment. \\
Scrum Master & Phan Anh Khang & Điều phối Planning/Daily/Review/Retro, theo dõi blocker và cập nhật hồ sơ. \\
Tech Lead / Schema Owner & Phan Anh Khang & Kiến trúc, schema, hợp đồng và tích hợp. PR của Khang vẫn cần người khác review. \\
Development Team & Khang, Khải, Hoàng, Kiệt, Huy & Phân tích, thiết kế, code, kiểm thử, review chéo và bàn giao. \\
\bottomrule
\end{tabularx}
\end{table}

\section{Bảng phân công công việc}
\label{sec:assignment}
Phân công theo mảng chính kết hợp chia đầu việc của từng luồng. Một
module có thể được nhiều người phối hợp chỉnh sửa; người được giao chịu
trách nhiệm đầu ra không nhất thiết là tác giả của mọi dòng code.
Bảng~\ref{tab:assignment} tổng hợp trách nhiệm, không quy đổi thành tỷ lệ
đóng góp hoặc xác nhận tất cả đầu việc đã hoàn thành.

{\small\setstretch{1.1}
\begin{longtable}{@{}L{0.22\textwidth}L{0.33\textwidth}L{0.38\textwidth}@{}}
\caption{Phân công theo thành viên Nhóm 7}\label{tab:assignment}\\
\toprule\textbf{Thành viên / MSSV} & \textbf{Mảng phụ trách chính} & \textbf{Đầu ra và phối hợp} \\\midrule\endfirsthead
\toprule\textbf{Thành viên / MSSV} & \textbf{Mảng phụ trách chính} & \textbf{Đầu ra và phối hợp} \\\midrule\endhead
\bottomrule\endfoot
Phan Anh Khang\newline 3124410145 & Nền tảng, schema/hợp đồng, Orders, tích hợp và Scrum. & Scaffold FE/BE, env/build/Compose; tạo đơn, snapshot, transaction, history/detail API; thiết kế state machine; phối hợp inventory/payment/E2E và hồ sơ Sprint. \\
Đoàn Quang Khải\newline 3124410138 & Frontend chi tiết, giỏ, checkout và màn hình sau thanh toán. & UI tùy chọn, cart store/identity/persist/tests; nối tạo đơn và payment; confirmation/history/detail; review các phần menu/auth/đầu ra Scrum theo thẻ. \\
Nguyễn Huy Hoàng\newline 3124410105 & Catalog, Inventory, dataset và seed. & Menu/detail API, menu UI và trạng thái; dữ liệu sản phẩm/size/topping; công thức và nguyên liệu; reserve/consume/release/concurrency; migration/seed demo. \\
Huỳnh Thái Kiệt\newline 3124410172 & Payments, Promotions và runtime state machine. & Mock gateway, payment outcome/retry/idempotency; coupon validation/lifecycle; phối hợp tác động PAID; full lifecycle core và extension staff queue theo kế hoạch. \\
Đỗ Mạnh Huy\newline 3124410113 & Auth, Security, Loyalty và tài khoản demo. & Register/login, bcrypt/JWT, session/logout/401; role/ownership; owner nghiệp vụ điểm thưởng, phối hợp PAID với Kiệt; hướng dẫn kiểm tra bảo mật và demo. \\
\end{longtable}}

\subsection{Phối hợp giữa các mảng}
\begin{table}[H]
\centering\small
\caption{Điểm phối hợp chính trong luồng nghiệp vụ}
\label{tab:collaboration}
\begin{tabularx}{\textwidth}{@{}L{0.22\textwidth}Y@{}}
\toprule\textbf{Luồng} & \textbf{Phân chia đầu việc} \\\midrule
Menu và chi tiết & Hoàng: API/dữ liệu/menu; Khải: UI tùy chọn; Khang: schema và hợp đồng. \\
Giỏ và tạo đơn & Khải: giỏ/request; Khang: Orders; Huy: JWT/principal; Hoàng: nguyên liệu; Kiệt: coupon. \\
Thanh toán và retry & Kiệt: payment/gateway; Khải: checkout; Khang: tích hợp/E2E; Hoàng: reservation/expiry; Huy: phiên và quyền. \\
Điểm thưởng & Huy: owner nghiệp vụ; Kiệt: phối hợp tác động PAID trong code; Khang: reviewer theo phân công. \\
Lịch sử và xác nhận & Khang: API/dữ liệu; Khải: UI; Huy: phiên/ownership; Kiệt: payment/state; Hoàng: review UI. \\
Demo và bàn giao & Khang: toàn stack/tài liệu; Hoàng: migration/seed/menu; Huy: tài khoản/bảo mật; các thành viên kiểm tra chéo theo thẻ. \\
\bottomrule
\end{tabularx}
\end{table}

Theo bản phân công đối chiếu, 40 thẻ hiện hành gồm 36 thẻ kỹ thuật core,
3 thẻ Scrum và 1 extension. Số thẻ giao lần lượt là Khang 12, Hoàng 10,
Huy 6, Khải 6, Kiệt 6. Kích thước thẻ khác nhau nên không dùng số thẻ để
suy phần trăm công sức. Danh mục mã thẻ và reviewer được trình bày ở
Phụ lục~\ref{sec:task-assignment}.

\section{Product Backlog và kế hoạch Sprint}
Bảng sau giữ mã, Sprint và story point từ Product Backlog trong
repository~\cite{backlog}. Đây là 10 hạng mục tổng hợp theo đề; bảng
không thay thế danh mục thẻ chi tiết hoặc tự đổi lịch cam kết trên board.
\begin{table}[H]
\centering\small\setstretch{1.1}
\caption{Product Backlog tổng hợp theo tài liệu nhóm}\label{tab:backlog}
\begin{tabularx}{\textwidth}{@{}L{0.12\textwidth}L{0.09\textwidth}YL{0.06\textwidth}@{}}
\toprule\textbf{Mã} & \textbf{Sprint} & \textbf{Hạng mục / kết quả mong đợi} & \textbf{SP} \\\midrule

BL-01 & 1 & Khởi tạo monorepo, env mẫu, lint/test và README. & 3 \\
BL-02 & 1 & API danh sách sản phẩm, seed và contract test. & 3 \\
BL-03 & 1 & Menu tích hợp API, loading/empty/error. & 3 \\
BL-04 & 1 & Chi tiết và size/topping/đường/đá/note hợp lệ. & 5 \\
BL-05 & 2 & Giỏ thêm/sửa/xóa, identity, tổng và persist. & 5 \\
BL-06 & 2 & Tạo đơn server-side, giữ tồn, trạng thái PENDING. & 8 \\
BL-07 & 2 & Đăng ký/đăng nhập, hash, JWT và ownership. & 5 \\
BL-08 & 3 & Thanh toán giả lập, idempotency, PAID/FAILED. & 8 \\
BL-09 & 3 & Xác nhận, lịch sử, Docker và demo E2E. & 8 \\
BL-10 & 3 & State machine, concurrency, coupon, loyalty và tests. & 13 \\
\bottomrule
\end{tabularx}
\end{table}

Mục tiêu Sprint 1 là nền tảng và menu; Sprint 2 là giỏ/tài khoản/tạo đơn;
Sprint 3 là thanh toán và hoàn thiện nghiệp vụ. Board hiện được tổ chức
thành S1 tuần 1--2, S2 tuần 3--4 và S3 tuần 5. Một số công việc như giỏ
hoặc auth được lên thẻ ở mốc khác với bảng backlog tổng hợp. Báo cáo giữ
sự phân biệt giữa hai cấp kế hoạch, không tự suy ngày họp, ngày kết thúc
hoặc kết quả Sprint khi thiếu biên bản.

\section{Kết quả đạt được}
Đối chiếu source và hồ sơ kiểm chứng cho thấy luồng khách hàng menu,
giỏ, đăng nhập, tạo đơn, thanh toán giả lập và xác nhận đã có tích hợp.
Backend có tính giá và snapshot, cơ chế giữ nguyên liệu theo công thức,
xử lý coupon, payment attempt và loyalty ledger. Lịch sử/chi tiết có API
và giao diện gắn với tài khoản. Dự án có migration, seed và nhiều mức
kiểm thử để kiểm tra các quy tắc quan trọng.

\section{Hạn chế và hướng hoàn thiện}
Ưu tiên tiếp theo là sửa lịch sử PAID; hoàn thiện state machine và luật
hủy; đưa probe idempotency/loyalty vào regression dùng DB thật; chốt chính
sách điểm; hoàn thiện bàn giao Compose/seed/tài khoản demo và hồ sơ
Scrum. Giao diện quầy được phát triển như extension sau khi các quy tắc
core được kiểm tra đầy đủ.

Các bước phát triển sau có thể bao gồm tích hợp cổng thanh toán thật,
quy trình đối soát/hoàn tiền, đo tải theo NFR-05 và cải thiện trải nghiệm
trên nhiều kích thước màn hình. Các nội dung này là định hướng, không
được tính vào kết quả hiện tại của BrewLite.
